%% kapa: exam=1 round=35 year=2024 period=1 form=A subject=civil_law subject_ko=민법 seq=39 number=39
%% source: Q-Net 감정평가사 기출문제 (cst00309) artlSeq=5213229 file=raw/1008/5213229/1교시 문제지 원본.pdf
%% conversion: kordoc 4.24.51 -> tools/gen1.py (auto); review=proofread-sonnet; flags=none
\begin{question}{39}
법정지상권이 성립하는 경우를 모두 고른 것은? (특별한 사정은 없고, 다툼이 있으면 판례에 따름)
\begin{framedtext}
ㄱ. X토지에 저당권을 설정한 甲이 저당권자 乙의 동의를 얻어 Y건물을 신축하였으나 저당권 실행 경매에서 丙이 X토지의 소유권을 취득한 경우 \par
ㄴ. 甲소유의 X토지와 그 지상건물에 공동저당권이 설정된 후 지상건물을 철거하고 Y건물을 신축하였고 저당권의 실행으로 X토지의 소유자가 달라진 경우 \par
ㄷ. X토지를 소유하는 甲이 乙과 함께 그 지상에 Y건물을 신축ㆍ공유하던 중 X토지에 저당권을 설정하였고 저당권 실행 경매에서 丙이 X토지의 소유권을 취득한 경우 \par
\end{framedtext}
\begin{choices}
  \item ㄱ
  \item ㄷ
  \item ㄱ, ㄴ
  \item ㄴ, ㄷ
  \item ㄱ, ㄴ, ㄷ
\end{choices}
\end{question}
\answer{2}
